% =====================================================================
% V-E. Análisis cualitativo de fallos  (E5)
% Fuente: resultados/f9_resumen_fallos.csv y figuras f9_*
% =====================================================================
\subsection{Análisis cualitativo de fallos}
\label{sec:fallos}

Buscamos los fallos del SAR que explican los errores de test y de Lima. Para no quedarnos en la
inspección visual, cada tipo de fallo se acompaña de una prueba cuantitativa independiente,
calculada con capas auxiliares descargadas sobre la grilla exacta del recorte
(Tabla~\ref{tab:fallos}). Para las pruebas espaciales usamos la \emph{razón de enriquecimiento}:
el porcentaje de errores que cumple una condición dividido entre el porcentaje del área que la
cumple. Si es mayor que 1, la condición concentra errores más de lo que su área explicaría.

\begin{table}[t]
\centering
\caption{Tipos de fallo del SAR y su evidencia cuantitativa.}
\label{tab:fallos}
\setlength{\tabcolsep}{3pt}
\begin{tabular}{@{}>{\raggedright\arraybackslash}p{2.4cm}c>{\raggedright\arraybackslash}p{4.6cm}@{}}
\toprule
Fallo (recorte) & Error & Evidencia \\
\midrule
Sombra / incidencia rasante (Lima\_0104) & FP &
  76.5\,\% de los FP con $\theta_{\text{local}}>60^\circ$ frente a 23.0\,\% del área
  (razón 3.3) \\
Doble rebote / lodo urbano (Lima\_0506) & FN &
  $\Delta$VV (17/03 $-$ 05/03): $+1.22$\,dB en zona inundada frente a $+0.03$\,dB fuera \\
Agua no especular (Pakistan\_849790) & FN &
  VH de los FN $-18.1$\,dB, cerca del no agua ($-16.3$) y lejos del agua detectada
  ($-24.9$) \\
Agua no especular (Bolivia\_314919) & FN &
  VH de los FN $-22.0$\,dB, intermedio entre agua ($-25.4$) y no agua ($-15.0$) \\
Superficie seca y lisa (Pakistan\_70625) & FP &
  VH de los FP $-24.9$\,dB, tan oscuro como el agua real ($-23.4$) \\
\midrule
\emph{Agua permanente (Lima\_0506)} & \emph{FP} &
  \emph{Descartada:} 0\,\% de los FP con ocurrencia JRC $>50$\,\% \\
\bottomrule
\end{tabular}
\end{table}

\textbf{1) Sombra de radar e incidencia rasante (FP).} Con el relieve SRTM y la geometría de la
pasada descendente calculamos el ángulo de incidencia local,
$\cos\theta_{\text{local}}=\cos s\cos\theta+\sin s\sin\theta\cos(a-\phi)$, donde $s$ es la
pendiente, $a$ la orientación de la ladera, $\theta$ el ángulo de incidencia y
$\phi\approx103^\circ$ el azimut hacia el sensor. Las laderas que miran en sentido contrario al
sensor reciben el pulso casi rasante y devuelven poca energía, igual que el agua. Los FP de la
quebrada de Lima\_0104 se concentran ahí: su $\theta_{\text{local}}$ mediano es $66.9^\circ$,
frente a $47.2^\circ$ en todo el recorte (Fig.~\ref{fig:caso1}).

\textbf{2) Doble rebote y lodo en zona urbana (FN).} En Lima\_0506 la inundación del Rímac
atraviesa zona urbana, y el modelo omite el 100\,\% de los píxeles inundados. La diferencia
entre la imagen del evento y la de doce días antes muestra que esa franja se \emph{aclaró}
(Fig.~\ref{fig:caso2}): es la firma del doble rebote y del barro rugoso. Esta es la causa
principal del colapso en Lima: el agua estaba, pero no se veía oscura.

\textbf{3) Agua no especular (FN).} En Pakistan\_849790 se omite el 84\,\% del agua. El Sentinel-2
muestra agua de aspecto somero junto a vegetación (Fig.~\ref{fig:caso3}), pero su VH es casi igual al de la
tierra: el viento, la vegetación emergente o la poca profundidad rompen la reflexión
especular. Es el mecanismo detrás del peor evento de test (Pakistán, IoU 0.194). En Bolivia\_314919 el
efecto es más leve y se limita a los bordes y zonas de transición.

\textbf{4) Superficie seca y lisa (FP).} Es el error inverso. En Pakistan\_70625 las parcelas en
barbecho o recién preparadas son tan lisas como el agua para la longitud de onda de banda C, y
su VH es incluso más oscuro que el del agua real (Fig.~\ref{fig:caso4}). Con una sola fecha de
SAR estos FP son físicamente indistinguibles del agua. No descartamos que una parte sea riego
no etiquetado, lo que haría el error ambiguo entre modelo y etiqueta. El mismo mecanismo
explica la mancha oscura de Lima\_0506, que ya era oscura antes del evento.

\textbf{Hipótesis descartada.} Probamos si los FP de Lima eran agua permanente que EMSN038 no
etiqueta como inundación. Con la ocurrencia de agua de JRC (1984--2021), ningún FP cae donde
hubo agua más de la mitad del tiempo, así que la descartamos. Reportarla es parte del análisis:
el patrón no se explica por una diferencia de convención de la etiqueta.

\begin{figure*}[t]
\centering
\includegraphics[width=0.85\textwidth]{f9_caso1_sombra.png}
\caption{Caso 1 (Lima\_0104): VH, incidencia local, probabilidad de agua y errores. Los FP
(rojo) coinciden con las laderas de incidencia rasante.}
\label{fig:caso1}
\end{figure*}

\begin{figure*}[t]
\centering
\includegraphics[width=0.85\textwidth]{f9_caso2_urbano.png}
\caption{Caso 2 (Lima\_0506): VV antes y durante, diferencia $\Delta$VV con el contorno de
EMSN038 y errores. La zona inundada se aclara (rojo en $\Delta$VV) y se omite por completo.}
\label{fig:caso2}
\end{figure*}

\begin{figure*}[t]
\centering
\includegraphics[width=0.85\textwidth]{f9_caso3_Pakistan_849790.png}
\caption{Caso 3 (Pakistan\_849790): el agua visible en Sentinel-2 no es oscura en VH y se omite
(naranja).}
\label{fig:caso3}
\end{figure*}

\begin{figure*}[t]
\centering
\includegraphics[width=0.85\textwidth]{f9_caso4_arido.png}
\caption{Caso 4 (Pakistan\_70625): parcelas secas y lisas tan oscuras como el agua generan FP
(rojo).}
\label{fig:caso4}
\end{figure*}
